
 %!TEX root = main.tex

\addcontentsline{toc}{section}{Use of Generative AI Tools}

This report was produced with the assistance of a generative AI
language model (Anthropic Claude) used as a productivity and
analytical support tool. Specifically, AI assistance was used for
the following tasks:

\begin{itemize}[leftmargin=*,itemsep=3pt]

  \item \textbf{Code generation and debugging.}
  The Python scripts implementing the Xpress MIP solver
  (Phase~1, Phase~2, resilient baseline, and Suez scenario),
  the min-cut and warehouse vulnerability analyses, the
  pipeline orchestration (\texttt{main.py}), and the interactive
  HTML visualiser were developed with AI-assisted code generation,
  iterative debugging, and structural review.

  \item \textbf{Data quality diagnosis.}
  The identification of the arc-capacity infeasibilities
  (A008 and A011 below supplier supply volumes) and the
  supply-demand imbalance for product~B were surfaced through
  AI-assisted cross-verification of the instance data against
  solver outputs.

  \item \textbf{Implementation of the resilient baseline.}
  The resilient baseline analysis was implemented with AI
  assistance, including the construction of synthetic bypass
  arcs, the application of the pessimism factor ($\times 1.5$),
  and the computation of the resilience ratio. The underlying
  analytical choices --- which suppliers to connect, which arcs
  to pre-activate, and how to interpret the results --- were
  made by the authors.

  \item \textbf{LaTeX drafting and structure.}
  Section drafts, table formatting, mathematical notation,
  and the overall report structure were refined with AI
  assistance. All analytical content, interpretations, and
  conclusions were reviewed, validated, and approved by the authors.

  \item \textbf{Critical review.}
  The AI tool was used to perform a simulated critical evaluation
  of the report prior to submission, identifying factual errors
  (supplier names, arc counts), copy-paste inconsistencies in
  tables, and gaps in the strategic assessment relative to the
  assignment rubric.

\end{itemize}

All quantitative results reported in this document were produced
by running the solver scripts on the provided instance data and
verified by the authors against the generated Excel outputs.
The AI tool did not access the solver directly and did not
produce any numerical results autonomously. Responsibility for
the analytical choices, model formulation, and conclusions
rests entirely with the authors.